%%%%%%%%%%%%%%%%%%%%%%%%%%%%%%%%%%%%%%
%%% Math commands %%%
%%%%%%%%%%%%%%%%%%%%%%%%%%%%%%%%%%%%%%
%% Generally useful

% Vector or matrix (boldface)
% \newcommand\vm[1]{\boldsymbol{\mathrm{#1}}}
\newcommand\vm[1]{\bm{#1}}


% R^n
\newcommand\R[1][\empty]{\mathbb{R}^{#1}}

% C^n
\newcommand\C[1][\empty]{\mathbb{C}^{#1}}

% Identity
\newcommand\I[1][\empty]{\vm{I}_{#1}}

% Transpose
\gdef\T{^{\top}}

% Inverse
\gdef\inv{^{-1}}

% Moore-Penrose Pseudoinverse
\gdef\mpp{^{+}}

% Set of elements
\newcommand\set[1]{\left\{#1\right\}}

% Span of vectors (\span was taken, sadly :( )
\newcommand\spn[1]{\textnormal{span}\left(#1\right)}

% Range of matrix
\newcommand\ran[1]{\textnormal{range}\left(#1\right)}

% Orthogonal complement
\gdef\orth{^{\perp}}

% Better trigonometric functions if one wants


% \let\oldcos\cos
% \renewcommand\cos[1]{\oldcos{\left(#1\right)}}
% \newcommand\scos[1]{\oldcos{^2\left(#1\right)}}

% \let\oldsin\sin
% \renewcommand\sin[1]{\oldsin{\left(#1\right)}}
% \newcommand\ssin[1]{\oldsin{^2\left(#1\right)}}

% \let\oldtan\tan
% \renewcommand\tan[1]{\oldtan{\left(#1\right)}}
% \newcommand\stan[1]{\oldtan{^2\left(#1\right)}}

% Hadamard product (slightly unnecessary? less readable? more readable?)
\let\hprod\odot

% Hadamard power
\newcommand\hpow[1]{^{\circ#1}}

% Kronecker product
\let\kron\otimes

% Slightly more meaningful names for pos.def, neg.def notation

% Pos. def, "matrix greater than"
\let\mgt\succ

% Neg. def
\let\mlt\prec

% Pos. semidef.
\let\mgeq\succcurlyeq

% Neg. semidef
\let\mleq\preccurlyeq

% Better real and imaginary part if one wants

% \let\oldRe\Re
% \renewcommand\Re[1]{\oldRe\left(#1\right)}

% \let\oldIm\Im
% \renewcommand\Im[1]{\oldIm\left(#1\right)}

%% Slightly more specific useful stuff

% Vectors with optional index and derivatives
\renewmathcommand\v[2][\empty]{\vm{#2}_{#1}}
\newcommand\dv[2][\empty]{\dot{\vm{#2}}_{#1}}
\newcommand\ddv[2][\empty]{\ddot{\vm{#2}}_{#1}}

% Scalars with optional index and derivatives
\newcommand\s[2][\empty]{#2_{#1}}
\newcommand\ds[2][\empty]{\dot{#2}_{#1}}
\newcommand\dds[2][\empty]{\ddot{#2}_{#1}}

% Vectors tilde with optional index and derivatives

\newcommand\vt[2][\empty]{\tilde{\vm{#2}}_{#1}}
\newcommand\dvt[2][\empty]{\dot{\tilde{\vm{#2}}}_{#1}}
\newcommand\ddvt[2][\empty]{\ddot{\tilde{\vm{#2}}}_{#1}}

% Vectors hat with optional index and derivatives
\newcommand\vh[2][\empty]{\hat{\vm{#2}}_{#1}}
\newcommand\dvh[2][\empty]{\dot{\hat{\vm{#2}}}_{#1}}
\newcommand\ddvh[2][\empty]{\ddot{\hat{\vm{#2}}}_{#1}}

% Vectors bar with optional index and derivatives
\newcommand\vb[2][\empty]{\bar{\vm{#2}}_{#1}}
\newcommand\dvb[2][\empty]{\dot{\bar{\vm{#2}}}_{#1}}
\newcommand\ddvb[2][\empty]{\ddot{\bar{\vm{#2}}}_{#1}}

% Scalars tilde with optional index and derivatives
\newcommand\st[2][\empty]{\tilde{#2}_{#1}}
\newcommand\dst[2][\empty]{\dot{\tilde{#2}}_{#1}}
\newcommand\ddst[2][\empty]{\ddot{\tilde{#2}}_{#1}}

% Scalars hat with optional index and derivatives
\newcommand\sh[2][\empty]{\hat{#2}_{#1}}
\newcommand\dsh[2][\empty]{\dot{\hat{#2}}_{#1}}
\newcommand\ddsh[2][\empty]{\ddot{\hat{#2}}_{#1}}

% Scalars bar with optional index and derivatives
% note: the first of these commands overwrite the "\sb{}" command which can be used instead of _{} to yield subscript. But who uses that?
\renewcommand\sb[2][\empty]{\bar{#2}_{#1}}
\newcommand\dsb[2][\empty]{\dot{\bar{#2}}_{#1}}
\newcommand\ddsb[2][\empty]{\ddot{\bar{#2}}_{#1}}

\newcommand\conj[1]{\overline{#1}}
\newcommand\upto[2][0]{\set{#1\dots #2}}
\newcommand\ceil[1]{\left\lceil#1\right\rceil}
\newcommand\floor[1]{\left\lfloor#1\right\rfloor}

%maybe figure these two out
\newcommand\der[2][]{\frac{\mathrm{d} #1}{\mathrm{d} #2}}
\newcommand\pder[2][]{\frac{\partial#1}{\partial#2}}
\newcommand\lder{\mathcal{L}}

\let\atan\arctan
\let\vphi\varphi
\DeclareMathOperator{\diag}{diag}

% Useful math operators
\DeclareMathOperator{\Ad}{Ad}
\DeclareMathOperator{\ad}{ad}
\DeclareMathOperator{\sgn}{\mathrm{sgn}}
\DeclareMathOperator{\blkdiag}{blkdiag}
\DeclareMathOperator{\SE}{SE}
\DeclareMathOperator{\se}{\mathfrak{se}}
\DeclareMathOperator{\SO}{SO}
\DeclareMathOperator{\so}{\mathfrak{so}}
\DeclareMathOperator{\rank}{rank}


\newcommand{\norm}[1]{\left\lVert#1\right\rVert}
\newcommand{\abs}[1]{\left|#1\right|}
\newcommand{\absign}[1]{\lceil#1\rfloor} % Signed absolute value



\gdef\i{\ensuremath{\mathrm{i}}}

% \DeclareMathOperator{\uh}{^{\wedge}}
\def\uh{^{\wedge}}
% \DeclareMathOperator{\lh}{^{\vee}}
\def\lh{^{\vee}}

% \newcommand\abs[1]{\left|#1\right|}

\makeatletter
\renewenvironment{proof}[1][\relax]{\par
  \pushQED{\qed}%
  \normalfont \topsep6\p@\@plus6\p@\relax
  \trivlist
  \item[\hskip\labelsep\bfseries
    \ifx#1\relax \proofname\else#1\fi\@addpunct{.}]\ignorespaces
}{%
  \popQED\endtrivlist\@endpefalse
}
\makeatother

\renewcommand{\qedsymbol}{$\blacksquare$}

\newtheorem{lemma}{Lemma}
\newtheorem{theorem}{Theorem}
\newtheorem{remark}{Remark}
\newtheorem{assumption}{Assumption}
\newtheorem{definition}{Definition}
\newtheorem{corollary}{Corollary}
\newtheorem{proposition}{Proposition}
%\newtheorem{prop}{Proposition}


\makeatletter
\providecommand{\examplesymbol}{\ensuremath\blacktriangle}
\newenvironment{example}[1][\examplename]{\par
  \pushQED{\qed}%
  \normalfont \topsep6\p@\@plus6\p@\relax
  \begingroup
  \let\qedsymbol\examplesymbol
      \trivlist

  \item[\hskip\labelsep
        \itshape
    #1\@addpunct{.}]\ignorespaces
}{%
  \popQED\endtrivlist\@endpefalse
  \endgroup
}
\providecommand{\examplename}{Example}
\makeatother

\crefname{prop}{Proposition}{Propositions}
\crefname{lemma}{Lemma}{Lemmas}
\crefname{theorem}{Theorem}{Theorems}
\crefname{assumption}{Assumption}{Assumptions}
\crefname{definition}{Definition}{Definitions}
\crefname{remark}{Remark}{Remarks}
\crefname{equation}{}{}
\crefname{corollary}{Corollary}{Corollaries}